% Описание алгоритма по схеме AlgoWiki:
% полностью дискретная полулагранжева схема для функции цены игры «шофёр-убийца»,
% решаемая методом итераций по значению (реализация: pursuit_game/hji_solver.py).
% Компилировать XeLaTeX или LuaLaTeX (как и основной файл курсовой).
\documentclass[14pt,a4paper]{extarticle}

\usepackage{fontspec}
\IfFontExistsTF{Times New Roman}{\setmainfont{Times New Roman}}{\setmainfont{Liberation Serif}}
\usepackage[russian]{babel}
\usepackage[a4paper, left=30mm, right=15mm, top=20mm, bottom=20mm]{geometry}
\usepackage{setspace}
\usepackage{indentfirst}
\usepackage{amsmath,amssymb}
\usepackage{float}
\usepackage{booktabs}
\usepackage{array}
\usepackage{algorithm}
\usepackage{algpseudocode}
\usepackage{tikz}
\usetikzlibrary{arrows.meta,calc,positioning,shapes.misc,shapes.geometric}

\setstretch{1.2}
\setlength{\parindent}{1.25cm}
\setlength{\emergencystretch}{3em}

\floatname{algorithm}{Алгоритм}
\renewcommand{\algorithmicrequire}{\textbf{Вход:}}
\renewcommand{\algorithmicensure}{\textbf{Выход:}}
\renewcommand{\algorithmicfor}{\textbf{для}}
\renewcommand{\algorithmicdo}{\textbf{выполнять}}
\renewcommand{\algorithmicif}{\textbf{если}}
\renewcommand{\algorithmicthen}{\textbf{то}}
\renewcommand{\algorithmicelse}{\textbf{иначе}}
\renewcommand{\algorithmicend}{\textbf{конец}}
\renewcommand{\algorithmicreturn}{\textbf{вернуть}}

\newcommand{\sgn}{\operatorname{sgn}}
\newcommand{\Tmax}{T_{\max}}

\begin{document}

\begin{center}
	{\bfseries\large Полностью дискретная полулагранжева схема\\
	для функции цены игры <<шофёр--убийца>>\\
	(метод итераций по значению)}
\end{center}

\vspace{0.5cm}

\tableofcontents
\newpage

%=====================================================================
\section{Свойства и структура алгоритма}
%=====================================================================

\subsection{Общее описание алгоритма}

Алгоритм предназначен для приближённого вычисления \emph{функции цены}
$V(x,y)$ дифференциальной игры преследования <<шофёр--убийца>>
в редуцированном пространстве~\cite{isaacs}. Платой в игре является время до захвата,
поэтому $V(x,y)$~--- это время, за которое преследователь~$P$ (шофёр)
гарантированно захватывает убегающего~$E$ из относительного положения
$(x,y)$ при оптимальном поведении обоих игроков.

Функция цены является (вязкостным) решением стационарного уравнения
Гамильтона--Якоби--Айзекса (основного уравнения Айзекса). Аналитически
это уравнение решается методом характеристик лишь по частям, между
сингулярными поверхностями; численный алгоритм позволяет получить $V$
сразу во всей расчётной области, включая сингулярные поверхности и
барьер.

Алгоритм относится к классу \emph{полулагранжевых схем} (semi-Lagrangian
schemes) для уравнений Гамильтона--Якоби--Беллмана--Айзекса. Описываемый
вариант~--- \emph{полностью дискретная схема} Барди--Фальконе--Соравиа
\cite{bfs,ff}: дискретизируются и время (шаг $\tau$), и пространство
(равномерная сетка с шагом $h$), а значения функции цены между узлами
восстанавливаются билинейной интерполяцией. Полученное сеточное
уравнение имеет вид неподвижной точки нелинейного оператора и решается
\emph{методом итераций по значению} (value iteration) в варианте Якоби:
на каждой итерации новое приближение во всех узлах вычисляется только по
предыдущему приближению.

Основные особенности алгоритма:
\begin{itemize}
	\item схема опирается на дискретный принцип динамического
	программирования: цена в узле равна минимаксу по управлениям от
	<<времени шага плюс цены в точке, куда приводит шаг>>;
	\item множества управлений заменяются конечными наборами:
	$\varphi\in\{-1,0,1\}$ для шофёра и $|\Psi|$ равномерно
	распределённых направлений $\psi$ для убегающего;
	\item захват засчитывается, если отрезок шага пересекает круг
	захвата (а не только если конец шага попал внутрь круга);
	\item в области, из которой убегающий уходит от захвата, точная цена
	бесконечна; в реализации она заменяется ограничением $\Tmax$
	(<<потолком>>).
\end{itemize}

\subsection{Математическое описание алгоритма}

\paragraph{Непрерывная задача.}
Редуцированная динамика игры (система координат связана с шофёром):
\begin{equation}\label{eq:dyn}
	\begin{aligned}
		\dot x &= f_1(x,y,\varphi,\psi) = -\frac{w_1}{R}\,y\,\varphi + w_2\sin\psi,\\
		\dot y &= f_2(x,y,\varphi,\psi) = \frac{w_1}{R}\,x\,\varphi - w_1 + w_2\cos\psi,
	\end{aligned}
\end{equation}
где $w_1$~--- скорость шофёра, $w_2<w_1$~--- скорость убегающего,
$R$~--- минимальный радиус поворота автомобиля,
$\varphi\in[-1,1]$~--- управление шофёра (кривизна траектории),
$\psi\in[0,2\pi)$~--- управление убегающего (направление движения).
Обозначим $f=(f_1,f_2)$, $z=(x,y)$.

Терминальное множество (круг захвата радиуса $l$):
$\mathcal{T} = \{ z : x^2+y^2 \le l^2 \}$.
Функция цены удовлетворяет основному уравнению Айзекса
\begin{equation}\label{eq:hji}
	\begin{aligned}
		&\min_{\varphi\in[-1,1]}\ \max_{\psi\in[0,2\pi)}
		\bigl[\, V_x f_1(z,\varphi,\psi) + V_y f_2(z,\varphi,\psi) + 1 \,\bigr] = 0,
		&& z\notin\mathcal{T};\\
		&V(z)=0, && z\in\mathcal{T}.
	\end{aligned}
\end{equation}

\paragraph{Дискретизация.}
Расчётная область $\Omega=[-L,L]^2$ покрывается равномерной сеткой
\[
	z_{ij}=(x_i,y_j),\quad x_i=-L+ih,\quad y_j=-L+jh,\quad i,j=0,\dots,n-1,
\]
где $n=2L/h+1$~--- число узлов по каждой координате.
Множества управлений заменяются конечными:
\[
	\Phi=\{-1,0,1\},\qquad
	\Psi=\Bigl\{\psi_k=\frac{2\pi k}{m}\ :\ k=0,\dots,m-1\Bigr\}.
\]
Для узла $z$ и пары управлений $(\varphi,\psi)$ точка, в которую система
переходит за время $\tau$ (явный шаг Эйлера вдоль характеристики), равна
\begin{equation}\label{eq:foot}
	z^{+}(z,\varphi,\psi) = z + \tau\, f(z,\varphi,\psi).
\end{equation}
Пусть $\theta(z,\varphi,\psi)\in[0,1]$~--- наименьший параметр, при котором
точка $z+\theta\,(z^{+}-z)$ попадает на окружность захвата
(корень квадратного уравнения $|z+\theta(z^{+}-z)|^2=l^2$),
если отрезок $[z,z^{+}]$ пересекает $\mathcal{T}$.

Обозначим через $I[V](z)$ билинейную интерполяцию сеточной функции
$V=\{V_{ij}\}$ в точке $z$, доопределённую так: $I[V](z)=0$ при
$z\in\mathcal{T}$ и $I[V](z)=\Tmax$ при $z\notin\Omega$.

\paragraph{Сеточное уравнение.}
Введём стоимость перехода
\begin{equation}\label{eq:cand}
	Q[V](z,\varphi,\psi)=
	\begin{cases}
		\tau\,\theta(z,\varphi,\psi), & \text{если отрезок } [z,z^{+}] \text{ пересекает } \mathcal{T},\\[2pt]
		\tau + I[V]\bigl(z^{+}(z,\varphi,\psi)\bigr), & \text{иначе,}
	\end{cases}
\end{equation}
и оператор
\begin{equation}\label{eq:op}
	(\mathcal{S}V)_{ij} =
	\begin{cases}
		0, & z_{ij}\in\mathcal{T},\\
		\min\Bigl\{\Tmax,\ \min\limits_{\varphi\in\Phi}\ \max\limits_{\psi\in\Psi}\ Q[V](z_{ij},\varphi,\psi)\Bigr\}, & z_{ij}\notin\mathcal{T}.
	\end{cases}
\end{equation}
Приближённая функция цены~--- неподвижная точка оператора: $V=\mathcal{S}V$.

\paragraph{Итерации по значению.}
Начальное приближение: $V^{(0)}_{ij}=0$ при $z_{ij}\in\mathcal{T}$ и
$V^{(0)}_{ij}=\Tmax$ иначе. Итерации
\begin{equation}\label{eq:iter}
	V^{(k)}=\mathcal{S}V^{(k-1)},\qquad k=1,2,\dots
\end{equation}
продолжаются, пока невязка
$r_k=\max_{ij}\bigl|V^{(k)}_{ij}-V^{(k-1)}_{ij}\bigr|$ не станет меньше
заданной точности $\varepsilon$ или пока не будет достигнуто
максимальное число итераций $K_{\max}$.

\paragraph{Синтез управлений.}
По найденной сеточной функции цены оптимальные управления в точке $z$
определяются из основного уравнения:
\[
	\sin\psi^{*}=\frac{V_x}{\sqrt{V_x^2+V_y^2}},\quad
	\cos\psi^{*}=\frac{V_y}{\sqrt{V_x^2+V_y^2}},\qquad
	\varphi^{*}=-\sgn\bigl(x V_y - y V_x\bigr),
\]
где $V_x,V_y$ вычисляются центральными разностями с шагом $h$.
Синтез управлений не входит в вычислительное ядро и выполняется
однократно для каждого запроса.

\paragraph{Значения параметров в реализации.}
$w_1=120$, $w_2=100$, $R=120$, $l=35$, $L=400$, $h=20$ ($n=41$),
$\tau=0{,}04$, $m=|\Psi|=24$, $\Tmax=8$, $\varepsilon=10^{-4}$,
$K_{\max}=400$.

\subsection{Вычислительное ядро алгоритма}

Вычислительное ядро~--- применение оператора $\mathcal{S}$
(формула~\eqref{eq:op}) ко всей сетке, т.\,е. одна итерация
\eqref{eq:iter}. На неё приходится практически всё время работы
алгоритма. Для каждого узла вне круга захвата и каждой из
$|\Phi|\cdot|\Psi|=72$ пар управлений выполняются:
\begin{enumerate}
	\item вычисление правой части \eqref{eq:dyn} и точки $z^{+}$ по
	формуле~\eqref{eq:foot};
	\item проверка пересечения отрезка $[z,z^{+}]$ с кругом захвата
	(решение квадратного уравнения);
	\item при отсутствии пересечения~--- билинейная интерполяция $I[V]$ в
	точке $z^{+}$ (чтение четырёх соседних значений $V^{(k-1)}$);
\end{enumerate}
после чего вычисляется максимум по $\psi$ (24 значения) и минимум по
$\varphi$ (3 значения).

Важное свойство ядра: точки $z^{+}$, признаки пересечения с кругом
захвата, индексы и веса интерполяции \emph{не зависят от $V$} и
одинаковы на всех итерациях. Поэтому ядро допускает вторую форму: после
однократного предвычисления <<таблицы переходов>> итерация сводится к
операции вида <<разреженная матрица $\times$ вектор>> (по 4 ненулевых
веса на пару <<узел--управление>>) с последующей редукцией min--max.

\subsection{Макроструктура алгоритма}

Алгоритм удобно описывать в терминах следующих макроопераций:
\begin{itemize}
	\item[\textbf{M1.}] \emph{Шаг характеристики}: вычисление $z^{+}$ по
	\eqref{eq:foot} для пары управлений.
	\item[\textbf{M2.}] \emph{Проверка захвата на отрезке}: вычисление
	$\theta(z,\varphi,\psi)$ или установление того, что пересечения нет.
	\item[\textbf{M3.}] \emph{Билинейная интерполяция} $I[V](z^{+})$:
	поиск ячейки сетки и взвешенная сумма четырёх значений.
	\item[\textbf{M4.}] \emph{Редукция min--max}: максимум по $\psi$
	для каждого $\varphi$ и минимум по $\varphi$ (поиск максимума и
	минимума в массиве).
	\item[\textbf{M5.}] \emph{Вычисление невязки}: поиск максимального
	элемента массива $|V^{(k)}-V^{(k-1)}|$ (глобальная редукция).
\end{itemize}
Одна итерация: для каждого узла M1--M3 для всех 72 пар управлений, затем
M4; по всем узлам~--- M5.

В варианте с предвычислением макрооперации M1, M2 и поиск ячейки и
весов из M3 выполняются один раз до итераций (макрооперация
\textbf{M0}~--- построение графа переходов), а в итерациях остаются
только выборка четырёх значений, взвешенная сумма, M4 и M5. Такой выбор
макроопераций делает информационную структуру итерации полностью
регулярной.

\subsection{Схема реализации последовательного алгоритма}

\begin{algorithm}[H]
	\caption{Итерации по значению для полулагранжевой схемы (вариант Якоби)}
	\label{alg:vi}
	\begin{algorithmic}[1]
		\Require $w_1,w_2,R,l,L,h,\tau,\Tmax,\varepsilon,K_{\max}$, наборы $\Phi$, $\Psi$
		\Ensure сеточная функция $V_{ij}$, число итераций $k$, невязка $r$
		\State $V_{ij}\gets 0$ при $x_i^2+y_j^2\le l^2$, иначе $V_{ij}\gets\Tmax$
		\For{$k=1,\dots,K_{\max}$}
		\State $r\gets 0$
		\For{$j=0,\dots,n-1$; $i=0,\dots,n-1$}
		\If{$x_i^2+y_j^2\le l^2$} $W_{ij}\gets 0$; \textbf{перейти к следующему узлу}
		\EndIf
		\State $b\gets\Tmax$ \Comment{лучшее для шофёра}
		\For{$\varphi\in\Phi$}
		\State $q\gets 0$ \Comment{худшее для шофёра (лучшее для убегающего)}
		\For{$\psi\in\Psi$}
		\State $z^{+}\gets z_{ij}+\tau f(z_{ij},\varphi,\psi)$ \Comment{M1}
		\If{отрезок $[z_{ij},z^{+}]$ пересекает $\mathcal{T}$ в точке с параметром $\theta$} \Comment{M2}
		\State $c\gets\tau\theta$
		\Else
		\State $c\gets\tau+I[V](z^{+})$ \Comment{M3}
		\EndIf
		\State $q\gets\max(q,c)$ \Comment{M4}
		\EndFor
		\State $b\gets\min(b,q)$ \Comment{M4}
		\EndFor
		\State $W_{ij}\gets\min(\Tmax,b)$;\quad $r\gets\max\bigl(r,|W_{ij}-V_{ij}|\bigr)$ \Comment{M5}
		\EndFor
		\State $V\gets W$
		\If{$r<\varepsilon$} \textbf{завершить}
		\EndIf
		\EndFor
	\end{algorithmic}
\end{algorithm}

Для итераций Якоби необходимы два массива размера $n\times n$
($V^{(k-1)}$ и $V^{(k)}$). Вариант Гаусса--Зейделя (обновление
массива на месте) требует одного массива и обычно сходится за меньшее
число итераций, но, как показано в п.~\ref{sec:par}, заменяет
координатный параллелизм на скошенный. Вариант с предвычисленной
таблицей переходов (п.~1.4) уменьшает время итерации ценой
дополнительной памяти порядка $n^2\cdot|\Phi||\Psi|\cdot 4$ индексов и
весов.

\subsection{Последовательная сложность алгоритма}\label{sec:seq}

Пусть $N=n^2$~--- число узлов, $P=|\Phi|=3$, $M=|\Psi|=24$, $K$~--- число
выполненных итераций. Оценка дана для вычислительного ядра в терминах
операций над вещественными числами.

На одну пару <<узел--управление>> в прямой реализации приходится
\begin{itemize}
	\item 2 вычисления тригонометрических функций ($\sin\psi$, $\cos\psi$;
	при табулировании по $\psi$~--- 0);
	\item около 14 умножений и сложений в M1;
	\item около 15 операций в M2 и, если дискриминант неотрицателен, одно
	извлечение квадратного корня;
	\item около 20 операций в M3 (из них 2 деления), если пересечения нет;
	\item 1 сравнение в M4.
\end{itemize}
Обозначим через $c\approx 50$ общее число арифметических операций на
пару. Тогда сложность одной итерации
\[
	W_1 \approx N\cdot P\cdot M\cdot c = 72\,c\,N \approx 3{,}6\cdot10^{3}\,N,
\]
плюс $N(PM-1)$ сравнений редукции min--max и $N$ сравнений для невязки.
Полная последовательная сложность
\[
	W = K\cdot N\cdot P\cdot M\cdot c = O\bigl(K\,n^2\,|\Phi|\,|\Psi|\bigr).
\]
Для параметров реализации ($N=1681$, $K=275$) это $W\approx 1{,}7\cdot 10^{9}$
операций.

Число итераций $K$ не известно заранее (выход по точности). Для схемы
с потолком $\Tmax$ оно по порядку не меньше $\Tmax/\tau$ (столько шагов
нужно, чтобы информация от круга захвата дошла до узлов с ценой,
близкой к $\Tmax$); на практике $K=114\div 275$ (см. п.~2.3).

В варианте с предвычислением построение таблицы переходов стоит
$\approx 72\,c\,N$ операций один раз, а итерация~--- около 8 операций
(4 умножения, 3 сложения, выбор ветви) и 4 чтения из памяти на пару, т.\,е.
$\approx 72\cdot 8\,N$ операций.

\subsection{Информационный граф}

Вершины графа соответствуют вычислению значения
$V^{(k)}_{ij}$ (с внутренними подвершинами для каждой пары
$(\varphi,\psi)$ и вершинами редукции M4) и вершинам редукции невязки M5.
Естественная система координат~--- $(k,j,i)$: номер итерации и индексы
узла. Граф состоит из $K$ одинаково устроенных ярусов-итераций.

\emph{Дуги внутри итерации} отсутствуют: значение $V^{(k)}_{ij}$
зависит только от значений предыдущей итерации $V^{(k-1)}$. Внутри
вершины узла $(k,j,i)$ данные идут от 72 подвершин <<шаг + интерполяция>>
к сдваивающим деревьям максимума (по 24 значениям) и минимума (по 3
значениям).

\emph{Дуги между итерациями} идут от вершин $(k-1,j',i')$ к вершине
$(k,j,i)$, где $(i',j')$~--- углы ячеек, содержащих точки $z^{+}$ для
всех пар управлений (по 4 угла на пару, если отрезок не пересекает
круг захвата). Все точки $z^{+}$ лежат в круге радиуса
$\tau\max|f|$ вокруг узла, поэтому дуги \emph{локальны}: шаблон зависимостей
имеет радиус
\[
	\rho = \Bigl\lceil \frac{\tau\,\max_{\Omega}|f|}{h}\Bigr\rceil + 1,\qquad
	\max_{\Omega}|f| \le \frac{w_1}{R}\sqrt2\,L + w_1 + w_2 .
\]
Для параметров реализации $\tau\max|f|\approx 31$, $h=20$, т.\,е.
$\rho\le 3$ узла; около круга захвата смещение меньше одной ячейки.
Форма шаблона различна в разных узлах (зависит от $z$ через
вращательное слагаемое в \eqref{eq:dyn}), но не зависит от значений $V$,
т.\,е. граф полностью определяется параметрами задачи.

Кроме того, от всех вершин итерации $k$ идут дуги к дереву редукции
невязки M5, результат которой (условие выхода) управляет наличием
итерации $k+1$.

\begin{figure}[H]
	\centering
	\begin{tikzpicture}[scale=0.95,
		pt/.style={circle,fill=gray!55,inner sep=1.3pt},
		pr/.style={circle,fill=black,inner sep=1.9pt},
		ft/.style={cross out,draw=red!70!black,thick,inner sep=2pt}]
		% layer k-1 grid
		\foreach \i in {0,...,6} \foreach \j in {0,...,6} \node[pt] at (\i,\j) {};
		\draw[gray!40] (0,0) grid (6,6);
		\node[anchor=north] at (3,-0.35) {ярус $k-1$: значения $V^{(k-1)}$};
		% circle of reachable foot points
		\draw[dashed,blue!60!black] (3,3) circle (1.75);
		\node[blue!60!black,anchor=west,font=\small] at (4.6,4.8) {$|z^{+}-z|\le\tau\max|f|$};
		% foot points and their cell corners
		\foreach \fx/\fy/\cx/\cy in {4.35/3.6/4/3, 2.3/4.45/2/4, 1.55/2.4/1/2} {
			\fill[orange!25] (\cx,\cy) rectangle ++(1,1);
			\node[pr] at (\cx,\cy) {}; \node[pr] at (\cx+1,\cy) {};
			\node[pr] at (\cx,\cy+1) {}; \node[pr] at (\cx+1,\cy+1) {};
			\node[ft] at (\fx,\fy) {};
		}
		\draw[very thick] (3,3) circle (0.17);
		\node[anchor=north east,font=\small] at (3,3) {$z_{ij}$};
		% layer k node
		\begin{scope}[xshift=9.2cm,yshift=3cm]
			\node[draw,thick,minimum width=1.1cm,minimum height=0.8cm] (v) at (0,0) {$V^{(k)}_{ij}$};
			\node[anchor=north,align=center,font=\small] at (0,-0.6) {ярус $k$};
		\end{scope}
		\draw[-{Stealth},thick] (5,4) to[bend left=18] (v.west);
		\draw[-{Stealth},thick] (3,5) to[bend left=25] (v.north west);
		\draw[-{Stealth},thick] (2,2) to[bend right=30] (v.south west);
		\node[align=left,font=\small] at (9.2,0.9) {{\color{red!70!black}$\times$} --- точки $z^{+}(z_{ij},\varphi,\psi)$\\
			$\bullet$ --- узлы-предшественники};
	\end{tikzpicture}
	\caption{Шаблон зависимостей одной вершины информационного графа:
		значение $V^{(k)}_{ij}$ зависит от значений $V^{(k-1)}$ в углах ячеек,
		содержащих точки $z^{+}$ (показаны три пары управлений из 72)}
	\label{fig:stencil}
\end{figure}

\begin{figure}[H]
	\centering
	\begin{tikzpicture}[
		box/.style={draw,thick,minimum width=1.5cm,minimum height=0.75cm,font=\small},
		red/.style={draw,thick,fill=gray!15,minimum width=1.9cm,minimum height=0.75cm,font=\small}]
		\foreach \k/\lab in {0/k-1,1/k} {
			\begin{scope}[yshift=-\k*2.6cm]
				\node[font=\small,anchor=east] at (-0.2,0) {итерация $\lab$};
				\foreach \c in {0,1,2} \node[box] (b\k\c) at (\c*2.0+0.8,0) {блок \c};
				\node[box,minimum width=1.0cm] (bd\k) at (6.6,0) {\dots};
				\node[red] (r\k) at (9.3,0) {$\max r$ (M5)};
				\draw[-{Stealth}] (bd\k.east) -- (r\k.west);
			\end{scope}
		}
		\foreach \c in {0,1,2} \draw[-{Stealth},thick,blue!60!black] (b0\c.south) -- (b1\c.north);
		\draw[-{Stealth},thick,blue!60!black] (b00.south) -- (b11.north);
		\draw[-{Stealth},thick,blue!60!black] (b01.south) -- (b10.north);
		\draw[-{Stealth},thick,blue!60!black] (b01.south) -- (b12.north);
		\draw[-{Stealth},thick,blue!60!black] (b02.south) -- (b11.north);
		\draw[-{Stealth},thick,blue!60!black] (bd0.south) -- (bd1.north);
		\draw[-{Stealth},dashed] (r0.south) -- node[right,font=\footnotesize,align=left]{условие\\выхода} (r1.north);
	\end{tikzpicture}
	\caption{Макроструктура информационного графа: ярусы-итерации,
		внутри яруса узлы независимы; дуги между соседними ярусами связывают
		только соседние блоки узлов (локальный шаблон радиуса $\rho$);
		глобальная редукция невязки определяет, будет ли следующий ярус}
	\label{fig:layers}
\end{figure}

\subsection{Ресурс параллелизма алгоритма}\label{sec:par}

\paragraph{Внутри итерации (массовый координатный параллелизм).}
В варианте Якоби все $N$ значений $V^{(k)}_{ij}$ информационно
независимы: вершины одного яруса лежат на гиперплоскости $k=\mathrm{const}$.
Внутри вершины узла независимы ещё все $P\cdot M=72$ подвершины
<<шаг + интерполяция>>. Ширина яруса~--- $N\cdot P\cdot M$
($\approx 1{,}2\cdot 10^{5}$ для $n=41$).

\paragraph{Высота одного яруса.}
При неограниченном числе процессоров итерация выполняется за
\[
	D_1 \approx d_c + \lceil\log_2 M\rceil + \lceil\log_2 P\rceil + \lceil\log_2 N\rceil
	\approx d_c + 5 + 2 + 11
\]
шагов, где $d_c\approx 15$~--- длина критического пути вычислений M1--M3
для одной пары (с извлечением корня), а логарифмические слагаемые~---
сдваивающие деревья редукций M4 и M5. В варианте с предвычислением
$d_c\approx 3$ (умножение и дерево из трёх сложений).

\paragraph{Между итерациями (конечный параллелизм отсутствует).}
Итерации выполняются строго последовательно. Таким образом,
параллельная сложность алгоритма
\[
	D = K\cdot\bigl(d_c + \lceil\log_2 (P M)\rceil + \lceil\log_2 N\rceil\bigr)
	= O(K\log N)
\]
при последовательной сложности $O(K\,N\,P\,M)$. Ярусы сбалансированы
по типу операций: на каждом ярусе одинаковый набор
арифметических операций, затем сравнения в редукциях, число операций
на уровнях деревьев редукции убывает вдвое.

\paragraph{Вариант Гаусса--Зейделя.}
При обновлении на месте с обходом по строкам вершина $(k,j,i)$ зависит
от уже вычисленных на той же итерации узлов шаблона с меньшими
номерами. Координатного параллелизма внутри итерации нет, но остаётся
\emph{скошенный} параллелизм: независимы узлы на <<фронтах>>
$i+\alpha j=\mathrm{const}$ с наклоном $\alpha$, определяемым радиусом
шаблона $\rho$. Распространённая альтернатива~--- раскраска узлов в
$(\rho+1)^2$ цветов, при которой узлы одного цвета обновляются
параллельно.

\paragraph{Иерархия параллелизма.}
Уровни (сверху вниз): узлы сетки (массовый, $N$) $\to$ управления
$\varphi$ (конечный, 3) $\to$ направления $\psi$ (массовый по $M$) $\to$
операции внутри M1--M3 (конечный). Верхний уровень естественно
использовать для распределения между процессами и нитями, нижние~--- для
векторизации.

\subsection{Входные и выходные данные алгоритма}

\textbf{Входные данные:} скаляры $w_1,w_2,R,l$ (параметры игры), $L,h,\tau,m$
(параметры дискретизации), $\Tmax,\varepsilon,K_{\max}$ (параметры
итераций). Всего $O(1)$ величин. Ограничения:
$w_1>w_2\ge 0$, $R>0$, $0<l<L$, $h>0$, $\tau>0$, $m\ge 4$.

\textbf{Выходные данные:} плотный массив $V_{ij}$ размера $n\times n$
вещественных чисел с диапазоном значений $[0,\Tmax]$ (значение $\Tmax$
означает, что за время $\Tmax$ захват не гарантирован), а также число
итераций, итоговая невязка и признак сходимости.

\textbf{Объём:} $N=n^2$ чисел на выходе; рабочая память~--- два массива
по $N$ чисел (Якоби) или дополнительно таблица переходов
$\approx N\cdot P\cdot M\cdot(4~\text{индекса}+4~\text{веса}+\text{флаги})$,
что в реализации на NumPy составляет около 5~КБ на узел.

Особенности данных: при параметрах реализации ($l/R\approx 0{,}29$)
конечные значения $V<\Tmax$ занимают лишь небольшую окрестность круга
захвата (например, при $h=20$ это 21 узел из 1681, включая 9 узлов круга захвата), остальная часть сетки
соответствует области ухода. Это согласуется с условием Айзекса: захват
из любой точки гарантирован лишь при
$l/R>\sqrt{1-\gamma^2}+\arcsin\gamma-1\approx 0{,}54$, где
$\gamma=w_2/w_1$~\cite{isaacs}.

\subsection{Свойства алгоритма}

\paragraph{Соотношение последовательной и параллельной сложности.}
$O(K N P M)$ против $O(K\log N)$: потенциальное ускорение линейно по
числу узлов и управлений. Последовательная часть~--- цепочка из $K$
итераций, которую нельзя распараллелить; уменьшать её можно только
уменьшением $K$ (метод Гаусса--Зейделя, ускоряющие обходы,
преобразование Кружкова, см. ниже).

\paragraph{Вычислительная мощность.}
Если считать входом и выходом одной итерации массивы $V^{(k-1)}$ и
$V^{(k)}$ ($2N$ чисел), вычислительная мощность прямой реализации равна
$72\,c\,N/(2N)\approx 1{,}8\cdot 10^{3}$ операций на элемент данных~---
ядро ограничено скоростью вычислений. В варианте с предвычислением на
каждую пару приходится $\approx 8$ операций на $\approx 70$ байт таблицы
переходов, т.\,е. итерация ограничена пропускной способностью памяти.
Для алгоритма в целом (вход $O(1)$, выход $N$) мощность
$\approx K\cdot 72\,c\sim 10^{6}$.

\paragraph{Устойчивость и сходимость.}
Веса билинейной интерполяции неотрицательны и в сумме равны единице, а
операции $\min$ и $\max$ монотонны, поэтому оператор $\mathcal{S}$
\emph{монотонен}: $U\le V\Rightarrow\mathcal{S}U\le\mathcal{S}V$. Отсюда
следует дискретный принцип максимума: значения итераций остаются в
$[0,\Tmax]$ при любом $\tau$, т.\,е. схема безусловно устойчива и не
требует условия Куранта. Начальное приближение $V^{(0)}=\Tmax$ вне
круга захвата является верхним решением, поэтому последовательность
$V^{(k)}$ монотонно не возрастает и сходится.
Сходимость дискретной цены к вязкостному решению \eqref{eq:hji} при
$h,\tau\to 0$ доказана для схем этого типа \cite{bfs,bcd}; погрешность
оценивается величиной порядка $O(\sqrt{\tau}+h/\tau)$, т.\,е. для
сходимости необходимо $h/\tau\to 0$: шаг сетки должен уменьшаться
быстрее шага по времени. При фиксированном $\tau$ измельчение сетки
не гарантирует уменьшения ошибки (это видно и по результатам п.~2.3).
Применение преобразования Кружкова $v=1-e^{-V}$ делает оператор
сжимающим с коэффициентом $e^{-\tau}$ и позволяет обойтись без
искусственного ограничения $\Tmax$ (область ухода получает $v=1$).

\paragraph{Детерминированность.}
Число итераций зависит от входных данных (выход по точности), поэтому
алгоритм в целом не детерминирован в смысле AlgoWiki. Однако структура
каждой итерации полностью детерминирована: ветвления в M2 и M3
(пересечение с кругом, выход за пределы сетки) зависят только от
геометрии и параметров, но не от значений $V$. Редукции M4 и M5 состоят
из операций $\min$ и $\max$, которые выполняются точно (без ошибок
округления), поэтому результат не зависит от порядка редукции, и
параллельная реализация воспроизводит последовательную побитово, если
сохранён порядок четырёх слагаемых в интерполяции. Расхождение между
двумя реализациями с разной группировкой слагаемых интерполяции составило
$1{,}8\cdot10^{-15}$ (п.~2.3).

\paragraph{Сбалансированность.}
Нагрузка на узлы почти одинакова. Исключения: узлы внутри круга захвата
(нулевая работа, $\approx\pi l^2/h^2$ узлов~--- менее 1\,\% сетки) и
пары управлений, для которых отрезок пересекает круг захвата
(интерполяция не нужна). Дисбаланс мал и при блочном распределении
узлов практически не проявляется.

\paragraph{Степень исхода и длинные дуги.}
Каждое значение $V^{(k-1)}_{ij}$ используется вершинами следующего
яруса из окрестности радиуса $\rho$~--- до нескольких сотен раз
(локальная рассылка). Длинных дуг, кроме рассылки скалярных параметров
и результата глобальной редукции невязки (условие выхода), нет. Это
благоприятно для размещения данных в кэш-памяти и для разбиения
области на подобласти.

\paragraph{Симметрия.}
Динамика \eqref{eq:dyn} инвариантна относительно замены
$(x,\varphi,\psi)\to(-x,-\varphi,-\psi)$, а наборы $\Phi$, $\Psi$
симметричны, поэтому $V(x,y)=V(-x,y)$. Вычисления можно вести только в
полуплоскости $x\ge 0$, что уменьшает объём работы и памяти вдвое.

%=====================================================================
\section{Программная реализация алгоритма}
%=====================================================================

\subsection{Особенности реализации последовательного алгоритма}

Базовая реализация (\texttt{pursuit\_game/hji\_solver.py},
класс \texttt{HJISolver}) написана на чистом Python и повторяет
алгоритм~\ref{alg:vi}: четыре вложенных цикла (строки, столбцы,
$\varphi$, $\psi$), функции \texttt{\_capture\_fraction} (M2) и
\texttt{\_interpolate} (M3), два массива для итераций Якоби, выход по
невязке $\varepsilon=10^{-4}$.

Возможности ускорения последовательной версии:
\begin{itemize}
	\item \emph{предвычисление таблицы переходов} (M0): так как $z^{+}$,
	признаки захвата, индексы и веса интерполяции не зависят от $V$, их
	можно вычислить один раз; итерация сводится к выборке и взвешенной
	сумме. Это компромисс <<память~--- время>>: около 5~КБ на узел;
	\item \emph{векторизация}: все пары <<узел--управление>> одной
	итерации обрабатываются как массивы (реализация на NumPy, п.~2.3);
	\item \emph{табулирование} $\sin\psi_k$, $\cos\psi_k$;
	\item \emph{использование симметрии} $V(x,y)=V(-x,y)$;
	\item \emph{метод Гаусса--Зейделя} с чередованием направлений обхода
	(fast sweeping): меньше итераций и один массив вместо двух;
	\item \emph{разрядность}: значения лежат в $[0,\Tmax]=[0,8]$, точность
	$\varepsilon=10^{-4}$, поэтому одинарной точности (float32,
	абсолютная погрешность округления $\sim 10^{-6}$) достаточно; это
	вдвое уменьшает объём таблицы переходов. Двойная точность нужна лишь
	при очень малых $\varepsilon$.
\end{itemize}

\subsection{Возможные способы и особенности параллельной реализации алгоритма}

\paragraph{Общая память (OpenMP, \texttt{multiprocessing}).}
Цикл по узлам (или по строкам сетки) внутри итерации распараллеливается
без синхронизаций, кроме барьера в конце итерации и редукции невязки
(\texttt{reduction(max:r)}). Внутренние циклы по $\psi$ подходят для
векторизации (SIMD).

\paragraph{Распределённая память (MPI).}
Сетка разбивается на полосы или прямоугольные блоки. Из-за локальности
шаблона каждому процессу нужны теневые грани (halo) шириной $\rho$ узлов.
На каждой итерации процесс:
\begin{itemize}
	\item выполняет $\approx K^{-1}W/p = 72\,c\,N/p$ операций;
	\item обменивается с соседями $2\rho n$ числами (полосовое разбиение) или
	$\approx 4\rho\, n/\sqrt p$ числами (блочное);
	\item участвует в одной глобальной редукции \texttt{MPI\_Allreduce(MAX)}
	невязки.
\end{itemize}
Отношение вычислений к обменам растёт как $n/\rho$, поэтому
масштабируемость хорошая уже при умеренных сетках. Результат редукции
по максимуму не зависит от порядка операций, так что результаты
воспроизводимы.

\paragraph{Графические ускорители (CUDA).}
Одна нить на узел (цикл по 72 управлениям внутри нити) или блок нитей
на узел с редукцией в разделяемой памяти. Таблицу переходов выгодно
хранить на ускорителе, тогда передача данных между хостом и ускорителем
нужна лишь в начале и в конце. Значения $V$ удобно читать через
текстурную память с аппаратной билинейной интерполяцией.

\paragraph{Вариант Гаусса--Зейделя.}
Параллельная реализация возможна через раскраску узлов или
волновые фронты (скошенный параллелизм); при этом число итераций
меньше, но ресурс параллелизма внутри итерации ниже.

\subsection{Результаты прогонов}

Платформа: виртуальная машина, Intel Xeon 2{,}1~ГГц, 2 вычислительных
ядра (vCPU), Python~3.13, NumPy~2.5.3. Параметры задачи~--- как в п.~1.2,
меняется только шаг сетки $h$. Сравниваются три реализации:
\begin{itemize}
	\item[A.] базовая последовательная (\texttt{HJISolver}, чистый Python);
	\item[B.] векторизованная (NumPy) с предвычисленной таблицей переходов;
	\item[C.] параллельная (\texttt{multiprocessing}, разбиение по строкам,
	ядро как в A).
\end{itemize}
Реализации A и B дают одинаковые числа итераций и невязки при всех $h$.
Максимальное расхождение значений $V$ при $h=40$ равно $1{,}8\cdot10^{-15}$.

\begin{table}[H]
	\centering
	\caption{Последовательные реализации A и B}
	\label{tab:seq}
	\small
	\begin{tabular}{@{}rrrrrrrr@{}}
		\toprule
		$h$ & $n\times n$ & $K$ & A, с & A, мс/итер. & B: табл., с & B: итер., с & B, мс/итер. \\
		\midrule
		40 & $21\times21$ & 152 & 6{,}60  & 43{,}4  & 0{,}03 & 0{,}09 & 0{,}58 \\
		20 & $41\times41$ & 275 & 49{,}75 & 180{,}9 & 0{,}10 & 0{,}56 & 2{,}04 \\
		10 & $81\times81$ & 164 & 112{,}95 & 688{,}7 & 0{,}16 & 1{,}93 & 11{,}78 \\
		5  & $161\times161$ & 114 & --- & --- & 0{,}52 & 6{,}09 & 53{,}42 \\
		\bottomrule
	\end{tabular}
\end{table}

Время итерации обеих реализаций растёт пропорционально числу узлов $N$, что
согласуется с оценкой п.~\ref{sec:seq}. Реализация B быстрее A в
$\approx 75$ раз при $h=20$ и в $\approx 54$ раза при $h=10$ (с учётом
построения таблицы). Таблица переходов занимает 8{,}7~МБ при $h=20$,
33{,}8~МБ при $h=10$ и 133{,}5~МБ при $h=5$.

\begin{table}[H]
	\centering
	\caption{Параллельная реализация C, $h=20$, 40 итераций}
	\label{tab:par}
	\begin{tabular}{@{}rrrr@{}}
		\toprule
		Процессов $p$ & с/итерацию & Ускорение & Эффективность \\
		\midrule
		1 & 0{,}182 & 1{,}00 & 1{,}00 \\
		2 & 0{,}123 & 1{,}48 & 0{,}74 \\
		\bottomrule
	\end{tabular}
\end{table}

Неполная эффективность при $p=2$ объясняется накладными расходами на
передачу массива $V$ процессам на каждой итерации (в \texttt{multiprocessing}
он сериализуется) и тем, что ядра виртуальной машины могут быть общими с
другими задачами. Для более крупных сеток доля этих расходов уменьшается.

\begin{table}[H]
	\centering
	\caption{Качество решения при измельчении сетки ($\tau=0{,}04$ фиксировано)}
	\label{tab:quality}
	\begin{tabular}{@{}rrrr@{}}
		\toprule
		$h$ & Узлов с $V<\Tmax-\tau$ & \multicolumn{2}{c}{Расстояние до барьера Айзекса} \\
		\cmidrule(l){3-4}
		 & & среднее & максимальное \\
		\midrule
		40 & 3  & 17{,}4 & 22{,}1 \\
		20 & 21 & 9{,}5  & 10{,}8 \\
		10 & 72 & 14{,}4 & 17{,}2 \\
		\bottomrule
	\end{tabular}
\end{table}

Расстояние от численной границы области с конечной ценой до
аналитического барьера~\cite{isaacs} уменьшается при переходе от $h=40$
к $h=20$, но не уменьшается при $h=10$. Это согласуется с оценкой
погрешности $O(\sqrt\tau+h/\tau)$: при фиксированном $\tau$ измельчение
сетки без уменьшения шага по времени не гарантирует сходимости, и для
расчётов на мелких сетках $\tau$ нужно уменьшать согласованно с $h$.

\subsection{Выводы для классов архитектур}

\begin{itemize}
	\item \emph{Многоядерные процессоры с общей памятью.} Итерации Якоби
	идеально подходят для OpenMP: независимые узлы, одна редукция по
	максимуму. Для сеток до $10^6$ узлов таблица переходов в одинарной
	точности помещается в оперативную память ($\sim 2{,}5$~ГБ), и
	производительность определяется пропускной способностью памяти;
	без таблицы ядро вычислительно ёмкое и хорошо векторизуется.
	\item \emph{Кластеры.} Разбиение области с теневыми гранями ширины $\rho$
	и одна глобальная редукция на итерацию дают хорошую масштабируемость;
	ограничивающий фактор~--- последовательная цепочка из $K$ итераций,
	поэтому для больших задач важно уменьшать $K$ (Гаусс--Зейдель
	по блокам, многосеточное начальное приближение, преобразование Кружкова).
	\item \emph{Графические ускорители.} Массовый параллелизм ширины
	$72N$, регулярный доступ к памяти и аппаратная интерполяция делают GPU
	наиболее подходящей платформой для мелких сеток и многомерных
	вариантов задачи (например, полной пятимерной постановки).
	\item \emph{Интерпретируемые языки.} Реализация на чистом Python на два
	порядка медленнее векторизованной; для расчётов на мелких сетках
	следует использовать NumPy, Numba или компилируемые языки.
\end{itemize}

%=====================================================================
\begin{thebibliography}{9}
	\bibitem{isaacs} Айзекс Р. Дифференциальные игры. --- М.: Мир, 1967. --- 480~с.
	\bibitem{bfs} Bardi M., Falcone M., Soravia P. Fully discrete schemes for the value function of pursuit-evasion games~// Advances in Dynamic Games and Applications. Annals of the International Society of Dynamic Games, vol.~1. --- Boston: Birkhäuser, 1994..
	\bibitem{ff} Falcone M., Ferretti R. Semi-Lagrangian Approximation Schemes for Linear and Hamilton--Jacobi Equations. --- Philadelphia: SIAM, 2014.
	\bibitem{bcd} Bardi M., Capuzzo-Dolcetta I. Optimal Control and Viscosity Solutions of Hamilton--Jacobi--Bellman Equations. --- Boston: Birkhäuser, 1997.
	\bibitem{pt} Patsko V.\,S., Turova V.\,L. Homicidal chauffeur game: history and modern studies~// Advances in Dynamic Games. Annals of the International Society of Dynamic Games, vol.~11. --- Boston: Birkhäuser, 2011..
	\bibitem{vv} Воеводин В.\,В., Воеводин Вл.\,В. Параллельные вычисления. --- СПб.: БХВ-Петербург, 2002. --- 608~с.
	\bibitem{algowiki} Открытая энциклопедия свойств алгоритмов AlgoWiki. Структура описания свойств алгоритмов. --- URL: https://algowiki-project.org
\end{thebibliography}

\end{document}
